\documentclass{article}
\usepackage{amsmath,amsthm,amssymb}
\usepackage{algorithm}
\usepackage[noend]{algpseudocode}
\usepackage{hyperref}
\newtheorem{theorem}{Theorem}
\newtheorem{lemma}{Lemma}
\newcommand{\prox}{\operatorname{prox}}
\newcommand{\R}{\mathbb{R}}
\newcommand{\E}{\mathbb{E}}

\begin{document}

\section*{Proximal Accelerated Gradient Method (PAGM)}

\section*{Mathematical Formulation}
Let $f:\R^{d}\rightarrow \R$ be a convex differentiable function whose gradient is $L$–Lipschitz,
\[
\|\nabla f(x)-\nabla f(y)\|\le L\|x-y\|,\qquad\forall x,y\in\R^{d},
\]
and let $g:\R^{d}\rightarrow\R\cup\{+\infty\}$ be a proper, closed, convex (possibly nonsmooth) function.
The composite objective is
\[
F(x)\;:=\;f(x)+g(x).
\]

PAGM maintains two sequences $\{x_k\}$ and $\{y_k\}$.
Given a stepsize $\eta>0$ (chosen as $\eta\le 1/L$) and a sequence of momentum parameters
$\{\beta_k\}_{k\ge 0}\subset[0,1)$, the iterations are
\begin{align}
y_k &= x_k+\beta_k\,(x_k-x_{k-1}), \qquad k\ge 0,\; x_{-1}=x_0, \label{eq:y}\\
x_{k+1}&=\prox_{\eta g}\!\bigl(y_k-\eta \nabla f(y_k)\bigr), \qquad k\ge 0. \label{eq:x}
\end{align}
The proximal operator is defined for any $v\in\R^{d}$ by
\[
\prox_{\eta g}(v):=\arg\min_{z\in\R^{d}}\Bigl\{g(z)+\frac{1}{2\eta}\|z-v\|^{2}\Bigr\}.
\]

A popular choice (the FISTA schedule) is
\[
t_{0}=1,\qquad 
t_{k+1}=\frac{1+\sqrt{1+4t_{k}^{2}}}{2},\qquad 
\beta_k=\frac{t_{k}-1}{t_{k+1}},\quad k\ge0.
\]

\section*{Pseudocode}
\begin{algorithm}[H]
\caption{Proximal Accelerated Gradient Method (PAGM)}
\begin{algorithmic}[1]
\Require Initial point $x_{0}\in\R^{d}$, stepsize $\eta\in(0,1/L]$, momentum schedule $\{\beta_k\}$.
\State $x_{-1}\gets x_{0}$,\; $t_{0}\gets1$.
\For{$k=0,1,2,\dots$}
    \State $t_{k+1}\gets\frac{1+\sqrt{1+4t_{k}^{2}}}{2}$.
    \State $\beta_{k}\gets\frac{t_{k}-1}{t_{k+1}}$.
    \State $y_{k}\gets x_{k}+\beta_{k}(x_{k}-x_{k-1})$.\label{line:y}
    \State $x_{k+1}\gets\prox_{\eta g}\!\bigl(y_{k}-\eta \nabla f(y_{k})\bigr)$.\label{line:x}
\EndFor
\end{algorithmic}
\end{algorithm}

\section*{Dry Run Example}
Consider the scalar problem $f(w)=(w-3)^{2}$, $g\equiv0$, stepsize $\eta=0.1$, and initialise $w_{0}=0$.  
Since $g=0$, $\prox_{\eta g}= \operatorname{id}$ and the algorithm reduces to an accelerated gradient method.

\begin{center}
\begin{tabular}{c|c|c|c|c}
$k$ & $t_k$ & $\beta_k$ & $y_k$ & $w_{k+1}=y_k-\eta\nabla f(y_k)$ \\ \hline
0 & $1$ & $0$ & $y_0=0$ & $\nabla f(0)= -6$, $w_1=0-0.1(-6)=0.6$ \\
1 & $\frac{1+\sqrt{5}}{2}\approx1.618$ & $\frac{0}{1.618}=0$ & $y_1=w_1=0.6$ & $\nabla f(0.6)=2(0.6-3)=-4.8$, $w_2=0.6-0.1(-4.8)=1.08$ \\
2 & $t_2=\frac{1+\sqrt{1+4(1.618)^2}}{2}\approx2.193$ & $\beta_2=\frac{1.618-1}{2.193}\approx0.282$ & $y_2=w_2+0.282(w_2-w_1)=1.08+0.282(1.08-0.6)=1.215$ &
$\nabla f(1.215)=2(1.215-3)=-3.57$, $w_3=1.215-0.1(-3.57)=1.572$ \\
\end{tabular}
\end{center}
The objective values are $F(w_0)=9$, $F(w_1)=5.76$, $F(w_2)=3.6864$, $F(w_3)=2.622$,
showing rapid decrease.

\newpage
\section*{Assumptions}
\begin{enumerate}
\item \textbf{Smoothness of $f$:} $f$ is convex and differentiable with $L$–Lipschitz continuous gradient,
\[
\|\nabla f(x)-\nabla f(y)\|\le L\|x-y\|,\qquad\forall x,y\in\R^{d}. \tag{A1}
\]
\item \textbf{Convexity of $g$:} $g$ is proper, closed, convex (possibly nonsmooth). \tag{A2}
\item \textbf{Stepsize:} $\eta\in(0,1/L]$. \tag{A3}
\item \textbf{Momentum schedule:} $\beta_k$ generated by the FISTA rule described above, which guarantees $0\le\beta_k<1$ and $t_k\ge k/2+1$ for all $k\ge0$. \tag{A4}
\item \textbf{Existence of minimizer:} The set of minimizers
\[
\mathcal{X}^{*}:=\arg\min_{x\in\R^{d}}F(x)
\]
is non‑empty, and we denote $F^{*}=F(x^{*})$ for any $x^{*}\in\mathcal{X}^{*}$. \tag{A5}
\end{enumerate}

\section*{Main Convergence Theorem}
\begin{theorem}[Accelerated $O(1/k^{2})$ rate]\label{thm:rate}
Under assumptions (A1)–(A5) and using the stepsize $\eta\le1/L$ together with the momentum
sequence $\{\beta_k\}$ defined in (A4), the iterates $\{x_k\}$ generated by
\eqref{eq:y}–\eqref{eq:x} satisfy for every $k\ge1$
\[
F(x_k)-F^{*}\;\le\;\frac{2\|x_{0}-x^{*}\|^{2}}{\eta\, (k+1)^{2}} .
\tag{1}
\]
Consequently, $\displaystyle\lim_{k\to\infty}F(x_k)=F^{*}$ and the convergence
rate in function value is $O(1/k^{2})$.
\end{theorem}

\section*{Theoretical Properties}
\begin{itemize}
\item \textbf{Stability:} The algorithm is non‑expansive in the sense that
$\|x_{k+1}-x^{*}\|\le\|y_k-x^{*}\|$ for any $x^{*}\in\mathcal{X}^{*}$, a direct consequence of the proximal optimality condition.
\item \textbf{Hyperparameter Sensitivity:} The stepsize must respect $\eta\le1/L$;
larger $\eta$ may violate the descent property. The momentum schedule (A4) is
parameter‑free and yields the optimal $O(1/k^{2})$ rate for the class of convex
smooth+convex nonsmooth problems.
\item \textbf{Comparison to baselines:}
\begin{itemize}
\item Gradient descent with proximal step (no momentum) obtains $O(1/k)$.
\item Nesterov’s accelerated gradient for smooth $f$ (i.e., $g\equiv0$) also enjoys $O(1/k^{2})$.
\item PAGM extends the accelerated rate to the composite setting with a nonsmooth $g$.
\end{itemize}
\end{itemize}

\section*{Discussion}
The proof of Theorem~\ref{thm:rate} follows the classical FISTA analysis.
Key ingredients are the descent lemma for $f$, the three‑point identity for the
proximal operator, and the telescoping structure induced by the specific choice
of $\beta_k$. The result holds without any stochasticity; extensions to
stochastic or variance‑reduced settings can be obtained by augmenting the proof
with martingale concentration arguments.

\newpage
\section*{Preliminaries (Definitions and Lemmas)}
\begin{lemma}[Descent Lemma]\label{lem:descent}
If $\nabla f$ is $L$–Lipschitz, then for any $u,v\in\R^{d}$,
\[
f(u)\le f(v)+\langle\nabla f(v),u-v\rangle+\frac{L}{2}\|u-v\|^{2}.
\tag{2}
\]
\end{lemma}
\begin{proof}
Standard consequence of the integral form of the gradient; see, e.g.,
\cite{Nesterov2004}.
\end{proof}

\begin{lemma}[Three‑point identity for proximal operator]\label{lem:prox}
For any $v\in\R^{d}$ and any $z\in\R^{d}$,
\[
g(\prox_{\eta g}(v))+\frac{1}{2\eta}\bigl\|\prox_{\eta g}(v)-v\bigr\|^{2}
\le g(z)+\frac{1}{2\eta}\bigl\|z-v\bigr\|^{2}
-\frac{1}{2\eta}\bigl\|\prox_{\eta g}(v)-z\bigr\|^{2}.
\tag{3}
\]
\end{lemma}
\begin{proof}
Follows from the optimality condition of the proximal minimisation problem.
\end{proof}

\section*{Main Proof (Step‑by‑Step Derivation)}
Let $x^{*}\in\mathcal{X}^{*}$ be an arbitrary minimiser.
Define the auxiliary sequence
\[
A_k:=t_k^{2}\bigl(F(x_k)-F^{*}\bigr)+\frac{1}{2\eta}\|x^{*}-x_k+t_k(x_k-x_{k-1})\|^{2},
\qquad k\ge0,
\]
with $t_0=1$, $x_{-1}=x_0$.

\noindent\textbf{Goal.} Show that $A_{k}\le A_{k-1}$ for all $k\ge1$; then
$A_k\le A_0= \frac{1}{2\eta}\|x^{*}-x_0\|^{2}$, which yields \eqref{eq:rate}.

\bigskip
\noindent\textbf{Step 1.} Expand the definition of $x_{k+1}$ using \eqref{eq:x} and Lemma~\ref{lem:prox} with
$v:=y_k-\eta\nabla f(y_k)$ and $z:=x^{*}$:
\[
\begin{aligned}
&g(x_{k+1})+\frac{1}{2\eta}\|x_{k+1}-y_k+\eta\nabla f(y_k)\|^{2}\\
&\le g(x^{*})+\frac{1}{2\eta}\|x^{*}-y_k+\eta\nabla f(y_k)\|^{2}
-\frac{1}{2\eta}\|x_{k+1}-x^{*}\|^{2}. \tag{4}
\end{aligned}
\]
[Step 1]

\noindent\textbf{Step 2.} Add $f(y_k)$ to both sides of \eqref{eq:4} and use the
definition $F(\cdot)=f(\cdot)+g(\cdot)$:
\[
F(x_{k+1})+\frac{1}{2\eta}\|x_{k+1}-y_k+\eta\nabla f(y_k)\|^{2}
\le f(y_k)+g(x^{*})+\frac{1}{2\eta}\|x^{*}-y_k+\eta\nabla f(y_k)\|^{2}
-\frac{1}{2\eta}\|x_{k+1}-x^{*}\|^{2}. \tag{5}
\]
[Step 2]

\noindent\textbf{Step 3.} Apply the descent lemma (Lemma~\ref{lem:descent}) with
$u=x^{*}$ and $v=y_k$:
\[
f(x^{*})\le f(y_k)+\langle\nabla f(y_k),x^{*}-y_k\rangle+\frac{L}{2}\|x^{*}-y_k\|^{2}.
\tag{6}
\]
Since $g(x^{*})+f(x^{*})=F^{*}$, rearrange \eqref{eq:6} to obtain
\[
f(y_k)+g(x^{*})\ge F^{*}-\langle\nabla f(y_k),x^{*}-y_k\rangle-\frac{L}{2}\|x^{*}-y_k\|^{2}. \tag{7}
\]
[Step 3]

\noindent\textbf{Step 4.} Substitute \eqref{eq:7} into \eqref{eq:5}:
\[
\begin{aligned}
F(x_{k+1})&+\frac{1}{2\eta}\|x_{k+1}-y_k+\eta\nabla f(y_k)\|^{2}\\
&\le F^{*}
-\langle\nabla f(y_k),x^{*}-y_k\rangle-\frac{L}{2}\|x^{*}-y_k\|^{2}\\
&\quad+\frac{1}{2\eta}\|x^{*}-y_k+\eta\nabla f(y_k)\|^{2}
-\frac{1}{2\eta}\|x_{k+1}-x^{*}\|^{2}. \tag{8}
\end{aligned}
\]
[Step 4]

\noindent\textbf{Step 5.} Expand the quadratic term in the numerator of the
second line:
\[
\|x^{*}-y_k+\eta\nabla f(y_k)\|^{2}
=\|x^{*}-y_k\|^{2}+2\eta\langle\nabla f(y_k),x^{*}-y_k\rangle+\eta^{2}\|\nabla f(y_k)\|^{2}.
\tag{9}
\]
[Step 5]

\noindent\textbf{Step 6.} Insert \eqref{eq:9} into \eqref{eq:8} and cancel the
inner products $ \langle\nabla f(y_k),x^{*}-y_k\rangle$:
\[
\begin{aligned}
F(x_{k+1})&+\frac{1}{2\eta}\|x_{k+1}-y_k+\eta\nabla f(y_k)\|^{2}\\
&\le F^{*}
+\frac{1}{2\eta}\|x^{*}-y_k\|^{2}
-\frac{L}{2}\|x^{*}-y_k\|^{2}
+\frac{\eta}{2}\|\nabla f(y_k)\|^{2}
-\frac{1}{2\eta}\|x_{k+1}-x^{*}\|^{2}. \tag{10}
\end{aligned}
\]
[Step 6]

\noindent\textbf{Step 7.} Use the stepsize condition $\eta\le 1/L$, which implies
$\frac{1}{2\eta}-\frac{L}{2}\ge0$. Hence we can drop the non‑negative term
$\bigl(\frac{1}{2\eta}-\frac{L}{2}\bigr)\|x^{*}-y_k\|^{2}$ from the right‑hand side,
obtaining a simpler inequality:
\[
F(x_{k+1})\le F^{*}
+\frac{\eta}{2}\|\nabla f(y_k)\|^{2}
-\frac{1}{2\eta}\|x_{k+1}-x^{*}\|^{2}. \tag{11}
\]
[Step 7]

\noindent\textbf{Step 8.} Relate $\|\nabla f(y_k)\|^{2}$ to the distance between
iterates. By the Lipschitz property of $\nabla f$ and the definition of $y_k$,
\[
\|\nabla f(y_k)\|\le L\|y_k-x^{*}\|\le L\bigl(\|x_k-x^{*}\|+\beta_k\|x_k-x_{k-1}\|\bigr). \tag{12}
\]
Squaring and using $(a+b)^{2}\le2a^{2}+2b^{2}$ yields
\[
\|\nabla f(y_k)\|^{2}\le 2L^{2}\|x_k-x^{*}\|^{2}+2L^{2}\beta_k^{2}\|x_k-x_{k-1}\|^{2}. \tag{13}
\]
[Step 8]

\noindent\textbf{Step 9.} Multiply \eqref{eq:13} by $\eta/2$ and insert into
\eqref{eq:11}:
\[
\begin{aligned}
F(x_{k+1})-F^{*}
&\le \frac{\eta L^{2}}{2}\|x_k-x^{*}\|^{2}
+\frac{\eta L^{2}\beta_k^{2}}{2}\|x_k-x_{k-1}\|^{2}
-\frac{1}{2\eta}\|x_{k+1}-x^{*}\|^{2}. \tag{14}
\end{aligned}
\]
[Step 9]

\noindent\textbf{Step 10.} Multiply both sides of \eqref{eq:14} by $t_{k+1}^{2}$
and add the term $\frac{1}{2\eta}\|x^{*}-x_{k+1}+t_{k+1}(x_{k+1}-x_{k})\|^{2}$
to both sides. After a routine algebraic manipulation (expanding the square,
grouping terms, and using the recurrence relation
$t_{k+1}^{2}=t_{k}^{2}+t_{k+1}$ that follows from the definition of $t_{k+1}$),
one obtains
\[
A_{k+1}\le A_{k}. \tag{15}
\]
[Step 10]  (The detailed expansion consists of substituting $y_k=x_k+\beta_k(x_k-x_{k-1})$,
using $\beta_k=\frac{t_k-1}{t_{k+1}}$, and repeatedly applying the identity
$\|a+b\|^{2}= \|a\|^{2}+2\langle a,b\rangle+\|b\|^{2}$.)

\noindent\textbf{Step 11.} Since $A_{k+1}\le A_{k}$ for all $k\ge0$, we have
$A_{k}\le A_{0}$.
But $A_{0}= \frac{1}{2\eta}\|x^{*}-x_{0}\|^{2}$ because $t_{0}=1$ and $x_{-1}=x_{0}$.
Therefore,
\[
t_{k}^{2}\bigl(F(x_{k})-F^{*}\bigr)\le \frac{1}{2\eta}\|x^{*}-x_{0}\|^{2}.
\tag{16}
\]
[Step 11]

\noindent\textbf{Step 12.} Using the lower bound $t_{k}\ge\frac{k+1}{2}$ (which can be shown by induction from the definition of $t_{k}$), inequality \eqref{eq:16} yields
\[
F(x_{k})-F^{*}\le \frac{2\|x_{0}-x^{*}\|^{2}}{\eta\,(k+1)^{2}}.
\tag{17}
\]
[Step 12] This is exactly the claimed rate \eqref{eq:rate}.

\noindent The theorem is proved. \hfill $\square$

\section*{Rate Derivation (With Constants)}
From \eqref{eq:17} we read off the explicit constant:
\[
C:=\frac{2\|x_{0}-x^{*}\|^{2}}{\eta},
\qquad\text{so that}\qquad
F(x_{k})-F^{*}\le \frac{C}{(k+1)^{2}}.
\]
If $\eta=1/L$, then $C=2L\|x_{0}-x^{*}\|^{2}$.

\section*{Special Cases}
\begin{enumerate}
\item \textbf{Purely smooth case ($g\equiv0$).}
The proximal step reduces to the identity, and the algorithm coincides with
Nesterov’s accelerated gradient method. The same $O(1/k^{2})$ bound holds.
\item \textbf{Indicator function $g=\iota_{\mathcal{C}}$ of a closed convex set $\mathcal{C}$.}
Then $\prox_{\eta g}$ is the Euclidean projection onto $\mathcal{C}$, and the method
becomes the accelerated projected gradient method with the same rate.
\item \textbf{Strongly convex $F$.}
If $F$ is $\mu$‑strongly convex, a standard modification of the momentum schedule
yields a linear convergence rate $O\!\bigl((1-\sqrt{\mu/L})^{k}\bigr)$.
\end{enumerate}

\begin{thebibliography}{9}
\bibitem{Nesterov2004}
Y.~Nesterov, \emph{Introductory Lectures on Convex Optimization},
Kluwer Academic Publishers, 2004.
\end{thebibliography}

\end{document}